\clearpage
\section{Tự luyện: từ đọc tham số đến chọn phương pháp}
\label{sec:practice}
Các bài dưới đây được biên soạn để luyện theo bài giảng. Khi làm, ghi cả
giả thiết và diễn giải kết quả, không chỉ viết một con số. Đáp án ở
Phụ lục~\ref{app:answers}; nên làm xong mỗi nhóm rồi mới đối chiếu.

\subsection{Nhóm A: tham số, tổ hợp và điều kiện}
\begin{enumerate}[label=\textbf{Bài \arabic*.},ref=\arabic*,series=practice]
\item Cho $\vmu=(10,20)'$,
$\Sigma=\begin{pmatrix}4&2\\2&9\end{pmatrix}$.
Hãy đọc trung bình, độ lệch chuẩn và tương quan; kiểm tra xác định dương;
tính $\Sigma^{-1}$. Vì sao chỉ biết hai phương sai chưa đủ để xác định
phân phối chuẩn hai chiều?

\item Với $\mathbf X\sim\N_3(\vmu,\Sigma)$ theo
\eqref{eq:running3d}, tìm phân phối của $(X_3,X_1)'$ và của $X_2$.
Viết ma trận chọn cho vectơ con. Giải thích vì sao không dùng công thức
có điều kiện ở đây.

\item Vẫn dùng \eqref{eq:running3d}, đặt $Y=X_1+2X_2-X_3$.
Tìm phân phối $Y$ bằng hai cách tính phương sai: nhân ma trận và khai triển
hiệp phương sai. Giải thích dấu của hạng chứa $\Cov(X_1,X_3)$.

\item Đặt $U=X_1+X_3$, $V=X_1-X_3$ trong cùng mô hình.
$U,V$ có độc lập không? $X_2,U$ có độc lập không? Nếu thay yêu cầu thứ
nhất bằng ``tìm $c$ sao cho $X_1+cX_3$ độc lập với $X_1-X_3$'', nghiệm là gì?

\item Cho $Z\sim\N(0,1)$ và $W=Z^2$.
Tính $\Cov(Z,W)$ và giải thích vì sao không suy ra độc lập.
Một người nói ``chỉ cần mỗi biến riêng chuẩn thì không tương quan tương
đương độc lập''. Hãy dùng phản ví dụ ở Phần~\ref{sec:linear} để bác bỏ.

\item Tìm phân phối $(X_1,X_2)'\mid X_3=4$ trong
\eqref{eq:running3d}. Viết đủ bốn khối hiệp phương sai cần dùng.
So sánh với phân phối biên và giải thích biến nào được cập nhật.
\end{enumerate}

\subsection{Nhóm B: phân phối mẫu và Hotelling}
\begin{enumerate}[label=\textbf{Bài \arabic*.},ref=\arabic*,resume=practice]
\item Mẫu iid $\N_2((10,20)',\begin{pmatrix}4&3\\3&9\end{pmatrix})$
có $n=25$. Nêu phân phối của $\Xbar$, phân phối của $24S$ và mối quan
hệ độc lập. So sánh độ lệch chuẩn của $X_{j1}$ và $\bar X_1$.
Nếu tổng thể không chuẩn nhưng vẫn có các moment hữu hạn, kết quả nào
về moment còn đúng và kết quả phân phối nào không còn chính xác?

\item Mẫu iid chuẩn có dữ liệu
\[
 \mathbf D=\begin{pmatrix}1&2\\2&1\\3&4\\4&3\\5&5\end{pmatrix}.
\]
Tính $\xbar,S,S^{-1}$; kiểm định $\vmu_0=(2,2)'$ mức 5\%.
Cho $F_{2,3;0.95}\approx9.5521$.
Tính thống kê có phân phối $F$ và viết kết luận bằng lời.
Viết log-hợp lý của mẫu; dùng phân rã tổng tích ngoài để tìm MLE trung
bình, rồi tìm MLE hiệp phương sai và so sánh với $S$.

\item Dùng mẫu Bài 8. Đặt
$C=\diag(100,0.1)$, $\mathbf d=(5,-2)'$.
Tính trung bình, hiệp phương sai và vectơ giả thuyết sau biến đổi.
Kiểm tra $T^2$ không đổi. Nếu bảng dữ liệu viết mỗi quan sát thành hàng,
phải nhân bảng với $C$ hay $C'$ trong công thức tổng quát?
Giải thích vì sao giá trị p và quyết định kiểm định cũng không đổi;
đẳng thức bất biến có cần giả thiết chuẩn không?
\end{enumerate}

\subsection{Nhóm C: khoảng đồng thời, hợp lý và mẫu lớn}
\begin{enumerate}[label=\textbf{Bài \arabic*.},ref=\arabic*,resume=practice]
\item Với số tóm tắt $n=40$ ở \eqref{eq:intervaldata}, cần báo cáo
$\mu_1,\mu_2,\mu_1-\mu_2,2\mu_1-\mu_2$.
Xác định $m$, tính ước lượng và sai số chuẩn của tổ hợp cuối.
Lập khoảng $T^2$ và Bonferroni cho tổ hợp đó.
Cho $t_{39;1-0.05/8}\approx2.6189$.
Có được dùng hệ số Bonferroni tính với $m=3$ cho cả bốn phát biểu không?
Xét riêng bài toán \emph{chỉ báo cáo hai trung bình}, dùng $m=2$ và
$t_{39;0.9875}\approx2.3313$. Dựng hình chữ nhật Bonferroni và kiểm tra
hai điểm $P=(10.70,18.95)'$, $R=(10.80,20.60)'$ trong chữ nhật và elip.
Nếu đổi $s_{12}$ từ 3 thành $-3$, miền nào đổi, và sai số chuẩn của
$2\mu_1-\mu_2$ tăng hay giảm?

\item Dùng Bài 8, tính $\widehat\Sigma$ và
$\widehat\Sigma_0$ dưới giả thuyết $(2,2)'$.
Từ định thức, tính $\lambda_W$, $\Lambda_{\mathrm{LR}}$ và $T^2$.
Vì sao Wilks' lambda không phải luôn là giá trị p?

\item Một mẫu iid từ tổng thể không chuẩn, có moment bậc hai hữu hạn,
cho $n=100,p=2$, $\xbar=(0.2,0.1)'$, $S=I_2$.
Với $H_0:\vmu=0$, tính kiểm định chi-bình phương xấp xỉ mức 5\% và
khoảng thành phần đồng thời chi-bình phương. Cho
$\chi^2_{2;0.95}=5.9915$, $z_{0.9875}=2.2414$.
Tính thêm khoảng Bonferroni cho hai thành phần.
Sau đó đánh giá hai phát biểu: ``$n=20,p=25$ vẫn dùng $S^{-1}$ bình thường'';
``$n=2000,p=1900$ đủ điều kiện xấp xỉ chỉ vì $n-p=100$''.
\end{enumerate}

\subsection{Tự đánh giá trước khi xem đáp án}
Bạn đã nắm ý chính khi có thể trả lời bằng lời:
\begin{itemize}
 \item Vì sao $\Sigma$ và $\Sigma/n$ xuất hiện trong hai bài toán khác nhau?
 \item Khi nào biết một biến khác giúp giảm phương sai có điều kiện?
 \item Tại sao $T^2$ có hệ số $n$ và phải chuyển đổi trước khi so với $F$?
 \item Miền elip khác hình chữ nhật từ khoảng thành phần như thế nào?
 \item Bonferroni bảo vệ danh sách nào, và ai quyết định số $m$?
 \item Chuẩn của trung bình mẫu có nghĩa là các quan sát gốc chuẩn không?
\end{itemize}
Nếu tính được số nhưng chưa giải thích được các câu này, hãy đọc lại phần
``Hiểu bằng lời'' tương ứng. Mục tiêu là chọn đúng công cụ cho câu hỏi,
không phải chỉ nhận ra hình dáng của một công thức.
